% GENERATED FILE. Edit canonical lessons, then run npm run generate:books.
% canonical-source-hash: fnv1a64:1224d194164c4054
% canonical-lessons: TA-C228-tattaccu-cey, TA-C228-molipeyar, TA-C228-rattu-cey, TA-C228-avacarappatu, TA-C228-neci

\chapter{To type, To translate, To cancel, To hurry, To love}
\label{ch:ta-to-type-to-translate-to-cancel-to-hurry-to-love}
\hlchaptermodality{228}

\begin{chapteropening}
\textbf{By the end of this chapter:} \emph{I can say to type, to translate, to cancel, to hurry and to love.}

Complete the last lesson of chapter 228: I can say to type, to translate, to cancel, to hurry and to love.
\end{chapteropening}

\section[\texorpdfstring{\ta{தட்டச்சு} \ta{செய்} (taṭṭaccu cey)}{தட்டச்சு செய் (taṭṭaccu cey)}]{\ta{தட்டச்சு} \ta{செய்} (taṭṭaccu cey) — to type}

\label{lesson:TA-C228-tattaccu-cey}



\begin{quote}
Before the new one: say the Tamil for to book in advance, then the Tamil for to apply.
\end{quote}

\subsection*{You'll want to know: \ta{தட்டச்சு} \ta{செய்}}
\textbf{\ta{தட்டச்சு} \ta{செய்}} — \emph{taṭṭaccu cey} — \textquotedblleft{}to type\textquotedblright{}.

\begin{grammarlens}[title={What you've built}]
One more word for this chapter.
\end{grammarlens}

\subsection*{Your turn}
\begin{itemize}
\raggedright
  \item \emph{Say it:} \emph{taṭṭaccu cey}
  \item \emph{Say it:} \emph{taṭṭaccu cey}, once more
  \item \emph{Say it:} say \emph{taṭṭaccu cey}
  \item \emph{Recall:} say the Tamil for to book in advance, then the Tamil for to apply, then say \emph{taṭṭaccu cey} again
\end{itemize}

\begin{tcolorbox}[breakable,colback=teal!4,colframe=teal!35!black,title={Before you move on}]
What does \ta{தட்டச்சு} \ta{செய்} mean? (\textquotedblleft{}To type\textquotedblright{}.) Say it once more.
\end{tcolorbox}
\section[\texorpdfstring{\ta{மொழிபெயர்} (moḻipeyar)}{மொழிபெயர் (moḻipeyar)}]{\ta{மொழிபெயர்} (moḻipeyar) — to translate}

\label{lesson:TA-C228-molipeyar}



\begin{quote}
Before the new one: say the Tamil for to apply, then the Tamil for to type.
\end{quote}

\subsection*{You'll want to know: \ta{மொழிபெயர்}}
\textbf{\ta{மொழிபெயர்}} — \emph{moḻipeyar} — \textquotedblleft{}to translate\textquotedblright{}.

\begin{grammarlens}[title={What you've built}]
One more word for this chapter.
\end{grammarlens}

\subsection*{Your turn}
\begin{itemize}
\raggedright
  \item \emph{Say it:} \emph{moḻipeyar}
  \item \emph{Say it:} \emph{moḻipeyar}, once more
  \item \emph{Say it:} say \emph{moḻipeyar}
  \item \emph{Recall:} say the Tamil for to apply, then the Tamil for to type, then say \emph{moḻipeyar} again
\end{itemize}

\begin{tcolorbox}[breakable,colback=teal!4,colframe=teal!35!black,title={Before you move on}]
What does \ta{மொழிபெயர்} mean? (\textquotedblleft{}To translate\textquotedblright{}.) Say it once more.
\end{tcolorbox}
\section[\texorpdfstring{\ta{ரத்து} \ta{செய்} (rattu cey)}{ரத்து செய் (rattu cey)}]{\ta{ரத்து} \ta{செய்} (rattu cey) — to cancel}

\label{lesson:TA-C228-rattu-cey}



\begin{quote}
Before the new one: say the Tamil for to type, then the Tamil for to translate.
\end{quote}

\subsection*{You'll want to know: \ta{ரத்து} \ta{செய்}}
\textbf{\ta{ரத்து} \ta{செய்}} — \emph{rattu cey} — \textquotedblleft{}to cancel\textquotedblright{}.

\begin{grammarlens}[title={What you've built}]
One more word for this chapter.
\end{grammarlens}

\subsection*{Your turn}
\begin{itemize}
\raggedright
  \item \emph{Say it:} \emph{rattu cey}
  \item \emph{Say it:} \emph{rattu cey}, once more
  \item \emph{Say it:} say \emph{rattu cey}
  \item \emph{Recall:} say the Tamil for to type, then the Tamil for to translate, then say \emph{rattu cey} again
\end{itemize}

\begin{tcolorbox}[breakable,colback=teal!4,colframe=teal!35!black,title={Before you move on}]
What does \ta{ரத்து} \ta{செய்} mean? (\textquotedblleft{}To cancel\textquotedblright{}.) Say it once more.
\end{tcolorbox}
\section[\texorpdfstring{\ta{அவசரப்படு} (avacarappaṭu)}{அவசரப்படு (avacarappaṭu)}]{\ta{அவசரப்படு} (avacarappaṭu) — to hurry}

\label{lesson:TA-C228-avacarappatu}



\begin{quote}
Before the new one: say the Tamil for to translate, then the Tamil for to cancel.
\end{quote}

\subsection*{You'll want to know: \ta{அவசரப்படு}}
\textbf{\ta{அவசரப்படு}} — \emph{avacarappaṭu} — \textquotedblleft{}to hurry\textquotedblright{}.

\begin{grammarlens}[title={What you've built}]
One more word for this chapter.
\end{grammarlens}

\subsection*{Your turn}
\begin{itemize}
\raggedright
  \item \emph{Say it:} \emph{avacarappaṭu}
  \item \emph{Say it:} \emph{avacarappaṭu}, once more
  \item \emph{Say it:} say \emph{avacarappaṭu}
  \item \emph{Recall:} say the Tamil for to translate, then the Tamil for to cancel, then say \emph{avacarappaṭu} again
\end{itemize}

\begin{tcolorbox}[breakable,colback=teal!4,colframe=teal!35!black,title={Before you move on}]
What does \ta{அவசரப்படு} mean? (\textquotedblleft{}To hurry\textquotedblright{}.) Say it once more.
\end{tcolorbox}
\section[\texorpdfstring{\ta{நேசி} (nēci)}{நேசி (nēci)}]{\ta{நேசி} (nēci) — to love}

\label{lesson:TA-C228-neci}



\begin{quote}
Before the new one: say the Tamil for to cancel, then the Tamil for to hurry.
\end{quote}

\subsection*{You'll want to know: \ta{நேசி}}
\textbf{\ta{நேசி}} — \emph{nēci} — \textquotedblleft{}to love\textquotedblright{}.

\begin{grammarlens}[title={What you've built}]
One more word for this chapter.
\end{grammarlens}

\subsection*{Your turn}
\begin{itemize}
\raggedright
  \item \emph{Say it:} \emph{nēci}
  \item \emph{Say it:} \emph{nēci}, once more
  \item \emph{Say it:} say \emph{nēci}
  \item \emph{Recall:} say the Tamil for to cancel, then the Tamil for to hurry, then say \emph{nēci} again
\end{itemize}

\begin{tcolorbox}[breakable,colback=teal!4,colframe=teal!35!black,title={Before you move on}]
What does \ta{நேசி} mean? (\textquotedblleft{}To love\textquotedblright{}.) Say it once more.
\end{tcolorbox}
